
% ============================================================
% Oppgave 1
% ============================================================
\hypertarget{oppgave1}{}
\section*{Oppgave 1}
\addcontentsline{toc}{section}{Oppgave 1}
Vi skal vise at konvolusjon er distributiv, altså at:
\[
x_1[n]*(x_2[n]+x_3[n])
=
x_1[n]*x_2[n]+x_1[n]*x_3[n]
\]
De tre sekvensene er:
\[
x_1[n]=\{4,-1,5\}
\]
\[
x_2[n]=\{6,-1,-1\}
\]
\[
x_3[n]=\{-4,3,2\}
\]
Først summeres $x_2[n]$ og $x_3[n]$:
\[
x_2[n]+x_3[n]
=
\{2,2,1\}
\]
Konvolusjonen blir dermed:
\[
x_1[n]*(x_2[n]+x_3[n])
=
\{4,-1,5\}*\{2,2,1\}
\]
Ved å utføre konvolusjonen får vi:
\[
\{8,6,12,9,5\}
\]
Vi beregner deretter konvolusjonene separat:
\[
x_1[n]*x_2[n]
=
\{24,-10,27,-4,-5\}
\]
og
\[
x_1[n]*x_3[n]
=
\{-16,16,-15,13,10\}
\]
Summen av disse blir:
\[
\{24,-10,27,-4,-5\}
+
\{-16,16,-15,13,10\}
=
\{8,6,12,9,5\}
\]
Begge metodene gir dermed samme resultat:
\[
\boxed{
x_1*(x_2+x_3)
=
x_1*x_2+x_1*x_3
}
\]
Konvolusjon er derfor distributiv.
\subsection*{Kontroll med lengre sekvenser}
Det samme testes med sekvensene:
\[
x_1[n]=\{4,-1,0,5\}
\]
\[
x_2[n]=\{6,-1,7,-1,1\}
\]
\[
x_3[n]=\{-4,3,2,1\}
\]
Siden $x_2$ og $x_3$ må ha samme lengde før de summeres, legges
det til en null bakerst i $x_3$:
\[
x_3[n]=\{-4,3,2,1,0\}
\]
I Octave beregnes dette som:
\begin{verbatim}
x1 = [4,-1,0,5];
x2 = [6,-1,7,-1,1];
x3 = [-4,3,2,1,0];

y1 = conv(x1, x2 + x3);
y2 = conv(x1,x2) + conv(x1,x3);
\end{verbatim}
Begge beregningene gir:
\[
y_1[n]
=
y_2[n]
=
\{8,6,34,1,14,44,0,5\}
\]
Dette viser at den distributive egenskapen også gjelder for de
lengre sekvensene.

